\title{Direct Preference Optimization: Your Language Model is Secretly a Reward Model}

\begin{abstract}
Although large-scale unsupervised language models (LMs) acquire extensive world knowledge and demonstrate certain reasoning abilities, precisely directing their outputs remains challenging because of the unsupervised methodology employed during their development. Prevailing techniques for achieving model controllability involve assembling human annotations on the relative merit of generated outputs, subsequently adapting the unsupervised LM through preference alignment—most commonly with reinforcement learning from human feedback (RLHF). Nonetheless, RLHF introduces significant complexity and instability by requiring the learning of a reward model to mirror human preferences, followed by reinforcement learning-based optimization of the LM to maximize the learned reward while preventing excessive deviation from the pretrained model. In this work, we propose a novel way to parameterize the reward model within the RLHF framework, which permits direct extraction of the optimal policy in closed form. This reformulation allows the conventional RLHF objective to be optimized using only a straightforward classification loss. The resulting method, termed \textit{Direct Preference Optimization} (DPO), achieves stable and effective alignment with modest computational demands, obviating the necessity for LM sampling during fine-tuning and minimizing hyperparameter adjustments. Empirical evaluations demonstrate that DPO aligns language models with human preferences on par with or better than previous techniques. Specifically, DPO surpasses PPO-based RLHF in guiding the sentiment of text outputs and achieves equal or superior performance in tasks like summarization and single-turn conversation, all while offering greater simplicity in implementation and training.
\end{abstract}

\section{Introduction}
Large-scale unsupervised language models (LMs) that are trained on extensive corpora have demonstrated a range of unexpected abilities~\citep{chowdhery2022palm, brown2020language, touvron2023llama, bubeck2023sparks}. Nevertheless, the data used for training such models is generated by humans with diverse intentions, levels of expertise, and priorities. Inevitably, some of these goals and skill sets are not ideal for imitation. For instance, although it is useful for an AI coding assistant to \textit{recognize} frequent programming errors so that it can fix them, we prefer that the model, when tasked with generating code, emulates the (possibly infrequent) high-caliber coding found in its data, rather than common errors. Similarly, while it may be beneficial for the model to be \textit{informed} about a widespread misconception held by half the population, we would certainly not want it to assert this misconception as fact in 50\% of the relevant prompts! In essence, the ability to \emph{select} the model’s target outputs and actions from among its expansive \textit{repertoire of knowledge and skills} is fundamental for constructing AI systems that are reliable, effective, and amenable to control \citep{ouyang2022training}. Although prevailing approaches predominantly rely on reinforcement learning (RL) to align LMs with human preferences, we demonstrate that the RL objective employed in these frameworks can in fact be exactly minimized using a straightforward binary cross-entropy loss, thereby streamlining the entire preference learning process.

Broadly speaking, current techniques instill target behaviors into language models by utilizing curated human preference datasets that embody the qualities users deem helpful and safe. This learning from preference occurs as a secondary phase, subsequent to large-scale unsupervised pre-training on massive textual corpora. The simplest strategy for learning from preferences is supervised fine-tuning based on exemplary human-generated outputs. Yet, the most effective methods are based on reinforcement learning from human or AI feedback (RLHF/RLAIF; \citep{christiano2017deep,bai2022constitutional}). In RLHF, a reward model is trained using human preferences, after which RL is used to adapt the language model's policy to favor outputs judged to be of high reward, while keeping deviations from the original model in check. Despite yielding models that exhibit notable performance in conversation and code generation, RLHF introduces substantial complexity relative to supervised learning, necessitating the training of multiple LMs and frequent policy sampling during the optimization loop, which considerably increases computational demands.

This work demonstrates a method for directly aligning a language model with human preferences, bypassing the need for explicit reward modeling or reinforcement learning strategies. We introduce \textit{{Direct Preference Optimization} (DPO)}, an approach that effectively targets the same optimization goal as standard RLHF methods (maximizing reward with a KL-divergence constraint), yet distinguishes itself through its ease of implementation and training. Conceptually, DPO functions by enhancing the relative log-likelihood of preferred responses compared to less favored ones, while incorporating a per-example, adaptive importance weight. This mechanism counteracts the issues of model collapse observed when using a naive probability ratio approach. Similar to prior work, DPO assumes an underlying theoretical preference model (e.g., the Bradley-Terry model; \cite{bradley1952rankanalysis}) for assessing the agreement between a reward function and observed preference data. In contrast to traditional pipelines, which employ the preference model to define a loss for reward model training before optimizing a policy for the reward, DPO reparameterizes the preference loss directly in terms of the policy. Given annotated preference data for model outputs, DPO enables direct policy optimization using a standard binary cross entropy loss, resulting in a policy that is optimal for an implicit reward function tailored to the collected preferences.

Our principal contribution, Direct Preference Optimization (DPO), provides a straightforward, RL-free alternative for preference-driven language model training. Empirical evaluations confirm that DPO matches or surpasses the performance of existing algorithms, such as PPO-based RLHF, across a range of tasks including sentiment control, summarization, and conversation, scaling to models with up to 6B parameters.

Large self-supervised language models, as their scale increases, acquire the capacity to solve certain problems with no task-specific training (zero-shot) \citep{radford2019language} or using a small number of task examples (few-shot prompts) \citep{gpt3,megatron,chowdhery2022palm}. Nonetheless, their outcomes on downstream applications and their adherence to user instructions are markedly enhanced through further training—specifically, fine-tuning—on datasets comprising instructions and human-produced completions \citep{mishra-etal-2022-cross,sanh2022multitask,chung2022scaling,thoppilan2022lamda}. This technique, known as `instruction-tuning,' equips LLMs with the ability to generalize to novel instructions outside the fine-tuning distribution, thus broadly improving model usability \citep{chung2022scaling}. In contrast to instruction-tuning, collecting \textit{relative} assessments from human evaluators regarding the quality of model outputs is typically more feasible than acquiring expert-generated demonstration data. Therefore, later research has turned to fine-tuning LLMs based on datasets capturing human preferences, which has led to better outcomes in tasks like translation \citep{kreutzer-etal-2018-reliability}, summarization \citep{stiennon2022learning,ziegler2020finetuning}, story-telling \citep{ziegler2020finetuning}, and instruction compliance \citep{ouyang2022training,ramamurthy2023is}. The prevalent approach in these settings involves first fitting a neural reward model aligned with observed preference data using frameworks such as the Bradley-Terry preference model \citep{bradley1952rankanalysis}, followed by adjusting the language model itself to optimize the learned reward, typically via reinforcement learning strategies such as REINFORCE \citep{williams1992reinforce}, proximal policy optimization (PPO; \cite{schulman2017proximal}), or their adaptations \citep{ramamurthy2023is}. Another complementary research direction makes use of LLMs trained to follow instructions with the aid of human feedback, prompting these models to generate supplementary artificial preference datasets targeting particular aspects like safety or harmlessness \citep{bai2022constitutional}. In this process, only minimal human guidance is required—specifically, textual rubrics for model evaluation. These methodologies draw together two prominent research threads: one that focuses on employing reinforcement learning for various language model objectives~\citep{Ranzato2015SequenceLT,paulus2018a,wu2018learning}, and another centered around generic frameworks for learning from human feedback \citep{christiano2017deep,kupcsik2018learning}. Although leveraging relative human preferences for model improvement offers promising avenues, applying reinforcement learning to fine-tune large language models remains a formidable practical barrier. The present work addresses this challenge by introducing a theoretically grounded alternative for optimizing with respect to relative human preferences without the use of RL.

Beyond language-focused settings, the study of learning policies from preferences has been explored in both bandit and reinforcement learning frameworks, leading to a variety of methodologies. In particular, contextual bandit learning that utilizes action preferences or rankings—rather than reward feedback—is termed the contextual dueling bandit (CDB; \cite{yue2012karmed,dudik2015contextual}). In cases where absolute rewards are absent, theoretical work on CDBs introduces the concept of a \textit{von Neumann winner}, defined as a policy achieving an expected victory rate of at least 50\% when compared against \textit{every} alternative policy \citep{dudik2015contextual}. A notable distinction is that CDB frameworks collect preference information in an online fashion, whereas learning from human preference typically relies on a static set of pre-collected preference-labeled action pairs \citep{yan2022human}. Analogously, \textit{preference-based RL} (PbRL) leverages binary preference data originating from an \textit{unknown} `scoring' mechanism as a proxy for reward signals \citep{BusaFekete2014,ruiz2023dueling}. There is a broad spectrum of PbRL algorithms, with several capable of utilizing off-policy preference datasets, though most methods involve first learning an explicit approximation of the underlying reward (or scoring) function and subsequently optimizing policies with respect to this estimate \citep{jain2013learning,BusaFekete2014,christiano2017deep,sadigh2017active,kupcsik2018learning}. In contrast, we propose a unified, single-stage approach that trains a policy to directly align with given preferences, bypassing the need for reward model estimation.

\section{Preliminaries}\label{section:prelims}

We begin by summarizing the RLHF framework detailed in \citeauthor{ziegler2020finetuning} and further extended in later works (\citep{stiennon2022learning, bai2022training, ouyang2022training}). The standard RLHF process generally consists of three stages: 1) supervised fine-tuning (SFT); 2) the collection of preference data along with reward model training; and 3) reinforcement learning optimization.

\textbf{SFT}: The RLHF procedure initiates by performing supervised fine-tuning on a pre-trained language model, using curated datasets tailored to the downstream application (such as dialogue generation, summarization, etc.), resulting in an adapted model $\pi^\text{SFT}$.

\textbf{Reward Modelling Phase}: During the second stage, the SFT model is prompted with an input $x$ to generate answer pairs $(y_1, y_2)\sim \pi^\text{SFT}(y \mid x)$. These pairs are then evaluated by human annotators, who select their preferred response, yielding labels $y_w\succ y_l \mid x$—where $y_w$ is the answer judged superior to $y_l$ for the given input. The source of these preferences is presumed to be an unknown underlying reward function $r^*(y, x)$, to which we do not have direct access. Preferences can be modeled in multiple ways, with the Bradley-Terry (BT) model \cite{bradley1952rankanalysis} being frequently used. More generally, ranking frameworks like Plackett-Luce (\citep{plackett1975analysis, luce2012individual}) can also be integrated if multi-way rankings are available. According to the BT formulation, the distribution $p^*$ over human preferences is expressed as follows:

\begin{equation}\label{eq:bradley-terry}
    p^*(y_1\succ y_2 \mid x)=\frac{\exp\left(r^*(x, y_1)\right)}{\exp\left(r^*(x, y_1)\right) + \exp\left(r^*(x, y_2)\right)}.
\end{equation}

Suppose we have a fixed dataset of comparisons $\mathcal{D}=\bigl\{x^{(i)}, y_w^{(i)}, y_l^{(i)}\bigr\}_{i=1}^N$ drawn according to $p^*$. We can introduce a parametrized reward function $r_{\phi}(x, y)$ and fit its parameters through maximum likelihood estimation. Treating this as a binary classification problem leads to the following negative log-likelihood objective:

\begin{equation}\label{eq:reward_model}
    \mathcal{L}_R(r_{\phi}, \mathcal{D}) = -\mathbb{E}_{(x, y_w, y_l)\sim \mathcal{D}}\bigl[\log \sigma(r_{\phi}(x, y_w)- r_{\phi}(x, y_l))\bigr]
\end{equation}

Here, $\sigma$ denotes the logistic sigmoid function. Within the framework of language models, $r_{\phi}(x, y)$ typically extends the SFT model $\pi^\text{SFT}(y \mid x)$ by adding a linear layer atop the last transformer block, which yields a scalar output corresponding to the reward estimate \cite{ziegler2020finetuning}. To minimize the variance of the estimated rewards, standard practice involves reward normalization so that $\mathbb{E}_{x,y\sim \mathcal{D}}\left[r_\phi(x, y)\right] = 0$ across all $x$.

\textbf{RL Fine-Tuning Phase}: In this stage, the acquired reward model is used as a feedback mechanism for the language model. As in previous research~\citep{jaques2017sequence, jaques2020human}, the problem is formulated as

\begin{equation}\label{eq:RL}
\max_{\pi_{\theta}}  \mathbb{E}_{x\sim \mathcal{D}, y\sim \pi_{\theta}(y \mid x)}\bigl[r_{\phi}(x, y)\bigr] - \beta\mathbb{D}_{\textrm{KL}}\bigl[\pi_{\theta}(y\mid x)\mid \mid \pi_\text{ref}(y\mid x)\bigr],
\end{equation}

where $\beta$ manages the allowable divergence from a reference policy $\pi_\text{ref}$, typically set as the initial SFT model $\pi^\text{SFT}$. The policy $\pi_\theta$ is usually initialized with the same SFT parameters. This constraint is essential to ensure that the updated policy does not stray significantly from the region where the reward model yields reliable outputs, while simultaneously helping to preserve output diversity and avoiding the collapse onto only a handful of high-reward responses. Because language modeling produces discrete outputs, the above objective is non-differentiable, so reinforcement learning algorithms are used for optimization. The prevalent method \citep{ziegler2020finetuning, stiennon2022learning, bai2022training, ouyang2022training} involves shaping the reward as ${r(x, y) = r_{\phi}(x, y) -\beta (\log \pi_{\theta}(y\mid x) - \log \pi_\text{ref}(y\mid x))}$ and maximizing it using PPO \cite{schulman2017proximal}.

\section{Direct Preference Optimization}\label{sec:DPO}

Recognizing the inherent complexities of applying RL algorithms to large-scale settings like language model fine-tuning, we seek an alternative, streamlined technique for policy optimization based directly on preference data. Diverging from the conventional RLHF paradigm, which entails learning a reward function and subsequently optimizing it via reinforcement learning, our method adopts a special parametrization of the reward model that admits a closed-form optimal policy, thus bypassing the need for an RL optimization loop. As elaborated below, the central innovation is an analytic connection between reward models and their associated optimal policies, allowing us to recast reward-model losses into losses on policies themselves. This reparameterization obviates the need for a separately-trained reward model, yet still grounds training in established frameworks for modeling preferences—such as the Bradley-Terry model. Consequently, the policy architecture implicitly fulfills both the roles of the generative model and the reward estimator.

\textbf{Derivation of the DPO objective.} Consider the reinforcement learning objective presented in Eq.~\ref{eq:RL}, equipped with a generic reward function $r$, consistent with earlier approaches. Building on the foundation laid by previous works~\citep{peters2007reinforcement, peng2019advantage, korbak2022reinforcement, go2023aligning}, it follows directly that the policy which maximizes the reward while being regularized by a KL term—as specified in Eq.~\ref{eq:RL}—assumes the following structure:

\begin{equation}\label{eq:op_policy}
    \pi_r(y\mid x) = \frac{1}{Z(x)}\pi_\text{ref}(y\mid x)\exp\left(\frac{1}{\beta}r(x, y)\right),
\end{equation}

with the partition function given by $Z(x) = \sum_{y}\pi_\text{ref}(y\mid x)\exp\left(\frac{1}{\beta}r(x, y)\right)$. Additional details are available in Appendix \ref{app:derivation1}. Although one might replace the unknown $r^*$ with its MLE-based estimator $r_{\phi}$, calculating $Z(x)$ remains computationally intensive \citep{korbak2022reinforcement, go2023aligning}, impeding the direct use of this policy formulation in applications. To circumvent this challenge, Eq.~\ref{eq:op_policy} can be algebraically rearranged to isolate the reward in terms of the optimal policy $\pi_r$, the baseline policy $\pi_\text{ref}$, and $Z(x)$. This is achieved by first taking the logarithm of both sides, yielding:

\begin{equation}\label{eq:main_eq}
    r(x,y) =\beta \log \frac{\pi_r(y\mid x)}{\pi_\text{ref}(y\mid x)} + \beta \log Z(x).
\end{equation}

Applying this transformation to both the ground-truth reward $r^*$ and its associated optimal policy $\pi^*$ is straightforward. Notably, within the Bradley-Terry framework, the only relevant term is the difference in rewards for two candidate responses; specifically, ${p^*(y_1 \succ y_2 \mid x) = \sigma(r^*(x, y_1) - r^*(x, y_2))}$. When the parametrization of Eq.~\ref{eq:main_eq} is substituted into this preference model (Eq.~\ref{eq:bradley-terry}) for $r^*(x,y)$, the partition function cancels, and the human preference likelihood can be reformulated using just $\pi^*$ and $\pi_\text{ref}$. Consequently, the optimal RLHF policy $\pi^*$, according to the Bradley-Terry assumption, obeys the following preference relationship:

\begin{equation}\label{eq:objective}
    p^*(y_1\succ y_2 \mid x)=\frac{1}{1 + \exp\left(\beta \log \frac{\pi^*(y_2\mid x)}{\pi_\text{ref}(y_2\mid x)} - \beta \log \frac{\pi^*(y_1\mid x)}{\pi_\text{ref}(y_1\mid x)}\right)}
\end{equation}

A full derivation is supplied in Appendix~\ref{app:derivation2}. Although Eq.~\ref{eq:objective} explicitly invokes the Bradley-Terry model, analogous derivations exist for broader models such as Plackett-Luce~\citep{plackett1975analysis, luce2012individual}, discussed in Appendix~\ref{app:plackett_luce_models}.

With the above, human preference probabilities can be directly related to the optimal policy rather than to an explicit reward parameterization, enabling the construction of a maximum likelihood objective for a parameterized policy $\pi_\theta$. Mirroring the paradigm employed in reward modeling (see Eq.~\ref{eq:reward_model}), we thus define the policy learning objective as:

\begin{equation}\label{eq:optimum_model}
    \mathcal{L}_\text{DPO}(\pi_{\theta}; \pi_\text{ref}) = -\mathbb{E}_{(x, y_w, y_l)\sim \mathcal{D}}\left[\log \sigma \left(\beta \log \frac{\pi_{\theta}(y_w\mid x)}{\pi_\text{ref}(y_w\mid x)} - \beta \log \frac

In this formulation, we approximate an implicit reward through an alternative parameterization, for which the optimal policy corresponds directly to $\pi_\theta$. Additionally, this methodology is mathematically akin to optimizing a reparametrized Bradley-Terry model, which provides us with desirable theoretical guarantees such as consistency, assuming appropriate conditions on the preference data distribution hold \cite{bong2022generalized}. Section~\ref{sec:theory} expands on the theoretical attributes of DPO and positions it relative to existing approaches.

\textbf{How does the DPO update operate?} To gain intuition into the inner workings of DPO, we examine the gradient of its loss function, $\mathcal{L}_\text{DPO}$. Specifically, the gradient with respect to the parameter set $\theta$ is given by:

\begin{multline*}\label{eq:gradient}
    \nabla_\theta \mathcal{L}_\text{DPO}(\pi_\theta;\pi_\text{ref}) = \\ -\beta\mathbb{E}_{(x, y_w, y_l) \sim \mathcal{D}} \bigg[\underbrace{\sigma(\hat{r}_\theta(x, y_l) - \hat{r}_\theta (x, y_w))}_\text{greater emphasis when reward estimate misranks} \bigg[\underbrace{\nabla_\theta\log \pi(y_w \mid x)}_\text{raise probability of $y_w$} - \underbrace{\nabla_\theta\log\pi(y_l \mid x)}_\text{lower probability of $y_l$}\bigg]\bigg],
\end{multline*}

Here, $\hat{r}_\theta(x, y) = \beta \log \frac{\pi_\theta(y \mid x)}{\pi_\text{ref}(y \mid x)}$ denotes the reward implicitly characterized by the language model $\pi_\theta$ and the baseline reference model $\pi_\text{ref}$ (see Section~\ref{sec:theory} for additional details). Conceptually, the gradient of $\mathcal{L}_\text{DPO}$ encourages the model to assign greater probability to responses labeled as preferred ($y_w$) while discouraging the model from favoring less desirable responses ($y_l$). Crucially, each example's influence is modulated by the degree to which the implicit reward model $\hat{r}_\theta$ erroneously assigns a higher value to the non-preferred outcome, with this contribution scaled by $\beta$, thereby reflecting both the degree of misranking and the effect of the KL regularization. Our empirical results highlight the necessity of this weighting; omitting the weighting term, as in a more straightforward variant, can lead to problematic degeneration in the language model (see Appendix Table~\ref{tab:unlikelihood_generations}).

\textbf{DPO framework.} The standard DPO workflow includes the following steps: 1) For each input prompt $x$, generate candidate completions $y_1, y_2 \sim \pi_\text{ref}(\cdot \mid x)$, and annotate them using human preference judgements to create the preference dataset $\mathcal{D} = \{x^{(i)}, y_w^{(i)}, y_l^{(i)}\}_{i=1}^N$; 2) Update the parameters of the language model $\pi_\theta$ to minimize $\mathcal{L}_\text{DPO}$, given the reference policy $\pi_\text{ref}$, the preference dataset $\mathcal{D}$, and the specified $\beta$.  
In applied settings, practitioners typically prefer to leverage existing preference datasets rather than acquiring new samples and collecting additional human labels. Since such datasets are generated using $\pi^\text{SFT}$, we set the initial reference policy as $\pi_\text{ref} = \pi^\text{SFT}$ when it is accessible. In situations where $\pi^\text{SFT}$ is unavailable, we obtain $\pi_\text{ref}$ by maximizing the likelihood over preferred responses ${(x, y_w)}$, specifically, by solving ${\pi_\text{ref} = \argmax_{\pi}\mathbb{E}_{x, y_w \sim \mathcal{D}}\left[\log \pi(y_w \mid x)\right]}$. This approach helps address the mismatch between the actual, inaccessible reference distribution and the practical $\pi_\text{ref}$ employed within DPO. Additional implementation details and choices of hyperparameters are provided in Appendix~\ref{app:implementation}.

\section{Theoretical Analysis of DPO} In this section, we deepen our understanding of DPO, present theoretical justification, and elucidate its benefits, particularly in comparison to the actor-critic family of algorithms utilized in RLHF, such as PPO~\cite{schulman2017proximal}.

\label{sec:theory}

\subsection{Your Language Model Is Secretly a Reward Model} DPO circumvents the need to construct an explicit reward model and eschews reinforcement learning for policy optimization, instead utilizing a unified maximum likelihood approach. Importantly, the DPO loss (see Eq. \ref{eq:main_eq}) aligns with a Bradley-Terry formulation where the reward is defined as $r^*(x, y) = \beta \log\frac{\pi^*_\theta(y \mid x)}{\pi_\text{ref}(y \mid x)}$. Optimization of the parameterized model $\pi_{\theta}$, in this context, mirrors the procedure of reward model fitting as described in Eq. \ref{eq:reward_model}, under an appropriate reparameterization. This section develops the formal groundwork underpinning this equivalence, demonstrates that it imposes no limitations on the space of possible reward models, and shows that the resulting approach can retrieve the optimal policy precisely. We commence by formalizing an equivalence relationship among reward functions.

\begin{definition}
We define two reward functions $r(x, y)$ and $r'(x, y)$ as equivalent if and only if there exists a function $f$ such that ${r(x, y)-r'(x, y) = f(x)}$.
\end{definition}

One can readily verify that this relation constitutes an equivalence relation, thus dividing the collection of reward functions into equivalence classes. This leads us to two central lemmas:

\begin{lemma}\label{lemma:same_prefrence} Within the Plackett-Luce framework—and by extension, the Bradley-Terry model—any pair of reward functions belonging to the same equivalence class generates identical preference distributions.
\end{lemma}

\begin{lemma}\label{lemma:same_policy}
    Any two reward functions in the same equivalence class will produce the same optimal policy under the specified constrained RL formulation.
\end{lemma}

The arguments supporting these claims are elementary, and thus we provide their details in Appendix \ref{app:lemma1}. The first lemma highlights a classical issue of model under-specification in the Plackett-Luce family \cite{plackett1975analysis}. As a result of this indeterminacy, additional identifiability constraints are typically introduced in order to guarantee well-defined MLE solutions for Eq.~\ref{eq:reward_model} \cite{bong2022generalized}. The second lemma asserts that, from a policy optimization standpoint, it suffices to identify any representative from the equivalence class, since all such reward functions induce the same optimal behavior. Based on these insights, we present the following theorem, with its proof located in Appendix~\ref{app:thm1}:

\begin{theorem}\label{thm:main}
    Under mild regularity conditions, every reward equivalence class compatible with the Plackett-Luce (including Bradley-Terry) models admits the reparameterized representation ${r(x, y) = \beta \log \frac{\pi(y\mid x)}{\pi_\text{ref}(y\mid x)}}$ for some policy model $\pi(y\mid x)$ and reference model $\pi_\text{ref}(y \mid x)$.
\end{theorem}

\begin{sproof}
    Given an arbitrary reward function $r(x, y)$, which specifies an associated optimal model $\pi_r(y \mid x)$ (see Eq.~\ref{eq:op_policy}), we aim to demonstrate that there exists an equivalent reward function, in the same class as $r$, that can be written in the prescribed reparameterized form. To this end, introduce the projection operator $f$ as
\begin{equation}
    f(r; \pi_\text{ref}, \beta)(x, y) = r(x, y) - \beta\log\sum_{y}\pi_\text{ref}(y\mid x)\exp\left(\frac{1}{\beta}r(x, y)\right)
\end{equation}
Here, $f$ effectively rescales the reward function by subtracting the logarithm of the partition function determined by $\pi_r$. Since this normalization term is dependent solely on $x$, $f(r; \pi_\text{ref}, \beta)(x, y)$ remains in the equivalence class of $r(x, y)$. Substituting $r$ on the right-hand side of Eq.~\ref{eq:main_eq} (valid for any reward function) gives $f(r; \pi_\text{ref}, \beta)(x, y) = \beta \log \frac{\pi_r(y\mid x)}{\pi_\text{ref}(y\mid x)}$. Therefore, $f$ yields a representative from the equivalence class of $r$ that adopts the desired log-ratio form, ensuring the generality of our reparameterized reward model.
\end{sproof}

Theorem~\ref{thm:main} can also be interpreted as identifying, within each class of equivalent reward functions, the specific reward function that the DPO reparameterization selects; namely, the reward function that fulfills the following condition:

\begin{equation}\label{eq:lag_p}
     \sum_{y}\underbrace{\pi_\text{ref}(y\mid x)\exp\left(\frac{1}{\beta}r(x, y)\right)}_{=\pi(y\mid x)\text{, using Thm.~\ref{thm:main} reparam.}} = 1,
\end{equation}

In other words, $\pi(y\mid x)$ defines a legitimate probability distribution (all probabilities are non-negative and sum to one). When considered alongside Eq.~\ref{eq:op_policy}, Eq.~\ref{eq:lag_p} emerges as the partition function for the optimal policy corresponding to the reward function $r(x, y)$. The principal contribution of the DPO algorithm lies in its ability to enforce specific constraints on the otherwise under-specified Plackett-Luce family of preference models (with a special focus on Bradley-Terry models), thus maintaining the expressive range of possible reward models while simultaneously ensuring that the optimal policy defined in Eq. \ref{eq:op_policy} is analytically tractable for any prompt $x$.

\subsection{Instability of Actor-Critic Algorithms} The framework we propose is also applicable for analyzing the sources of instability found in conventional actor-critic algorithms deployed in RLHF pipelines, such as PPO. Our focus is on the reinforcement learning fine-tuning stage as detailed in Section \ref{section:prelims}. There is a notable connection to the control as inference paradigm~\cite{levine2018reinforcement} in the context of the constrained RL formulation provided in \ref{eq:RL}. By considering a parameterized distribution $\pi_{\theta}(y\mid x)$, the goal is to minimize $\mathbb{D}_{\text{KL}}[\pi_{\theta}(y|x) \mid \mid \pi^*(y\mid x)]$, where $\pi^*$ is the optimal policy from Eq.~\ref{eq:optimum_model} and is derived from the reward function $r_{\phi}(y, x)$. After some algebraic manipulations, this minimization translates to the following optimization criterion:

\begin{equation}\label{eq:AC}
    \max_{\pi_{\theta}}\mathbb{E}_{\pi_{\theta}(y\mid x)}\bigg[\underbrace{r_{\phi}(x, y) -\beta\log\sum_{y}\pi_\text{ref}(y\mid x)\exp\left(\frac{1}{\beta}r_{\phi}(x, y)\right)}_{f(r_{\phi}, \pi_\text{ref}, \beta)} - \underbrace{\beta\log\frac{\pi_{\theta}(y\mid x)}{\pi_\text{ref}(y\mid x)}}_{\text{KL}}\bigg]
\end{equation}

The objective here aligns with that of previous studies \citep{ziegler2020finetuning, stiennon2022learning, bai2022training, ouyang2022training}, which also optimize the DPO-equivalent reward associated with the reward class $r_{\phi}$. In this context, the normalization term present in $f(r_{\phi}, \pi_\text{ref}, \beta)$ may be viewed as the soft value function corresponding to the reference policy $\pi_\text{ref}$. Although this normalization factor does not alter the set of optimal solutions, omitting it can lead to policy gradients with substantial variance, potentially destabilizing the learning dynamics. To address this, one option is to estimate the normalization via a learned value function; however, this approach poses its own optimization challenges. An alternative strategy employed in earlier literature is to use a baseline derived from human completions, effectively serving as a one-sample Monte Carlo estimate for the normalization. Distinctively, the DPO reparameterization provides a reward function formulation that is baseline-independent.

\section{Experiments}
Here, we systematically investigate the empirical performance of DPO when learning policies from preferences. Our initial experiments use a controlled text-generation framework to examine how effectively DPO balances reward maximization against KL-divergence from the reference policy, relative to widely used preference-based learning methods like PPO. Subsequently, we assess DPO’s scalability and capability on larger neural models and more challenging RLHF benchmarks, including tasks in summarization and conversational agents. Our findings suggest that DPO, with minimal hyperparameter adjustments, generally matches or surpasses the performance of strong baselines such as PPO-based RLHF, and also outperforms strategies that simply select the best of $N$ generated trajectories based on a learned reward model. Prior to presenting our empirical results, we outline our experimental methodology; supplementary information is provided in Appendix~\ref{app:exp_details}.

\textbf{Tasks.} Our empirical analysis investigates three distinct open-ended text generation tasks. Across all scenarios, the algorithms are trained to optimize a policy based on a preference dataset $\mathcal{D}=\bigl\{x^{(i)}, y_w^{(i)}, y_l^{(i)}\bigr\}_{i=1}^N$. For \textbf{controlled sentiment generation}, $x$ represents a partial movie review extracted from the IMDb corpus \cite{maas-EtAl:2011:ACL-HLT2011}, and the objective is to generate a continuation $y$ exhibiting positive sentiment. To ensure a standardized evaluation protocol, we synthetically construct preference pairs using outputs scored by a pre-trained sentiment classifier such that $p(\text{positive}\mid x,y_w)>p(\text{positive}\mid x,y_l)$. For supervised fine-tuning (SFT), GPT-2-large is adapted until training loss stabilizes, utilizing the training subset of the IMDB dataset (see further methodological specifics in App~\ref{app:sentiment_details}). In the case of \textbf{summarization}, $x$ corresponds to a Reddit forum post and the generation target $y$ is a concise summary encapsulating the primary arguments. Adopting established conventions, we employ the Reddit TL;DR summarization dataset \citep{volske-etal-2017-tl} paired with human-annotated preferences from \citeauthor{stiennon2022learning}. Here, the SFT baseline is trained on genuine forum post summaries,\footnote{\url{https://huggingface.co/CarperAI/openai_summarize_tldr_sft}} leveraging the TRLX \citep{leandro_von_werra_2023_7790115} platform for RLHF optimization. The human preference data was compiled by \citeauthor{stiennon2022learning} utilizing outputs from a comparably fine-tuned SFT model distinct from our primary baseline. Lastly, in the \textbf{single-turn dialogue} setting, $x$ is a single user prompt, which may range from technical queries about astrophysics to inquiries seeking personal guidance. Here, the policy is tasked with producing informative and engaging replies $y$, using the Anthropic Helpful and Harmless dialogue collection \citep{bai2022training}, comprised of 170k interactions between a person and an automated assistant. Each conversation concludes with two alternative responses produced by a (large, unspecified) language model and annotated according to human preference. In this context, there is no pre-existing SFT baseline; accordingly, we fine-tune a publicly available language model solely on the preferred continuations to construct the SFT model.

\textbf{Evaluation.} Our experimental framework utilizes two distinct evaluation strategies. To systematically assess how effectively each algorithm maximizes the constrained reward objective, we examine their performance in the controlled sentiment generation environment. Here, the reward-KL-divergence frontier is used as the metric, since the true reward function—implemented via a sentiment classifier—is fully observable. In practical scenarios where access to the true reward function is unavailable, we measure algorithm performance using their \textit{win rate} against a reference policy. This comparison leverages GPT-4, which acts as a stand-in for human judgments to evaluate summary quality in summarization and response helpfulness in the single-turn dialogue context. For summarization, the standard of comparison is the ground-truth summary from the test set; in dialogue, it is the test set’s preferred response. Notably, prior research indicates that language models can serve as more reliable automatic evaluators than traditional metrics \citep{Chen2023ExploringTU}. Nevertheless, we support our choice of GPT-4 for evaluation with a human user study, as described in Sec.~\ref{sec:human-judgments}. Our findings show that human-GPT-4 agreement is strong, frequently matching or surpassing inter-human annotator consistency.

\textbf{Methods.} Alongside DPO, our study benchmarks a variety of established methods for aligning language model outputs with human preferences. For the summarization task, we experiment with zero-shot prompting using \textbf{GPT-J} \citep{gpt-j}, and for dialogue, we adopt a 2-shot prompting paradigm with \textbf{Pythia-2.8B} \citep{biderman2023pythia}. We also include the \textbf{SFT} model, and \textbf{Preferred-FT}, which is fine-tuned in a supervised manner on selected completions $y_w$—derived from either the SFT model (for sentiment control and summarization) or a general language model (for dialogue). Another supervised variant is \textbf{Unlikelihood}~\citep{welleck2019neural}, which aims to boost the likelihood of $y_w$ while explicitly reducing that of $y_l$, regulated by an adjustable $\alpha\in[0,1]$ parameter governing the unlikelihood term. Further, we implement \textbf{PPO} \citep{schulman2017proximal} trained on a reward function distilled from preference data, and \textbf{PPO-GT}, an idealized approach that leverages the exact reward function available in sentiment-controlled experiments. For PPO-GT in sentiment control, we compare a standard implementation \cite{leandro_von_werra_2023_7790115} with an enhanced version that incorporates reward normalization and refined hyperparameters; these improvements are likewise used with PPO where rewards are learned. Lastly, we examine the \textbf{Best of $N$} method: sampling $N$ responses from either SFT (or Preferred-FT for dialogue) and selecting the one that scores highest under a learned reward function. Despite strong performance, this approach is computationally burdensome at inference since every test-time prompt necessitates generating $N$ completions.

\subsection{How well can DPO optimize the RLHF objective?}

In RLHF, the standard approach involves maximizing a reward function while simultaneously enforcing a KL-divergence constraint to prevent the learned policy from straying too far from a reference policy. As a result, any evaluation of different methods should jointly consider both the attained reward and the KL distance; obtaining marginally better rewards at the expense of a much greater KL is not always a beneficial outcome. Figure~\ref{fig:frontier-tldr-main} presents the trade-off curve between reward and KL for several methods within the sentiment classification framework. For each method, we conduct several training experiments, where each trial employs a distinct setting for the policy's conservativeness parameter (target KL $\in\{3,6,9,12\}$ for PPO, $\beta \in \{0.05,0.1,1,5\}$, $\alpha\in\{0.05,0.1,0.5,1\}$ for unlikelihood, and various random seeds for preferred-FT), resulting in a total of 22 experiments. At intervals of every 100 training steps, up to convergence, each model is assessed on a test suite of prompts. We record both the mean reward (according to the true reward function) and the mean sequence-level KL\footnote{Defined as the cumulative sum of KL-divergences across timesteps.} relative to the reference policy $\text{KL}\left(\pi\mid \mid \pi_\text{ref}\right)$. The outcomes indicate that DPO yields the most advantageous frontier, attaining the highest rewards with consistently low KL values. This finding is especially striking for several reasons. Firstly, although DPO and PPO target the same optimization objective, DPO demonstrates substantially higher efficiency; its reward/KL trade-off curve strictly dominates that of PPO. Secondly, DPO establishes a superior frontier compared to PPO, \emph{even when PPO is granted access to exact ground truth rewards} (PPO-GT).

\subsection{Can DPO scale to real preference datasets?}
\label{sec:dpo-real-datasets}
We subsequently assess the effectiveness of DPO fine-tuning on tasks such as summarization and single-turn dialogue. For summarization tasks, it is well established that automated metrics like ROUGE are not always strongly aligned with human judgment~\citep{stiennon2022learning}, and earlier research has shown that employing PPO for LM fine-tuning based on human feedback leads to higher quality summaries. In our analysis, we test various approaches by generating outputs on the TL;DR summarization dataset’s test split and calculating their average win rate in comparison to reference summaries from the test set. To ensure a comprehensive evaluation, we generate completions across a temperature range from 0.0 to 1.0, presenting the resulting win rates in Figure~\ref{fig:frontier-tldr-main} (right). DPO, PPO, and Preferred-FT are all applied as fine-tuning strategies to the identical GPT-J SFT model\footnote{\url{https://huggingface.co/CarperAI/openai_summarize_tldr_sft}}. Our results indicate that DPO attains a win rate near 61\% at temperature 0.0, surpassing PPO, which achieves about 57\% under its optimal setting (temperature 0.0). DPO further exhibits a higher peak win rate in comparison with the strongest $N$ baseline. Importantly, no significant tuning of DPO’s $\beta$ hyperparameter was performed, so these outcomes may be a conservative estimate of DPO’s full capability. Furthermore, DPO demonstrates much greater resilience to changes in sampling temperature, whereas PPO’s effectiveness deteriorates to the level of the base GPT-J model at elevated temperatures. Preferred-FT does not result in meaningful gains over the SFT baseline. Additionally, in a direct human evaluation discussed in Section~\ref{sec:human-judgments}, DPO outputs sampled at temperature 0.25 were preferred 58\% of the time over PPO outputs sampled at the same temperature.

For single-turn dialogue, we assess various approaches using the portion of the Anthropic HH dataset’s test split that contains a single exchange between human and assistant \citep{bai2022training}. For the GPT-4 evaluations, we use the preferred test completions as references to measure the win rate across different strategies. Since a standard SFT baseline is lacking for this specific task, our methodology involves initializing with the pre-trained Pythia-2.8B model. We subsequently apply Preferred-FT on the selected completions to train a reference model, ensuring that its completions are representative of the training distribution, followed by DPO training. Our comparisons include the top completion out of 128 Preferred-FT samples (as performance saturates beyond $N=128$; see Appendix Figure~\ref{fig:best-of-n}) and a 2-shot prompt version of the original Pythia-2.8B, finding that DPO matches or surpasses alternatives when evaluated at optimal temperatures for each method. Additionally, we examine an RLHF model, trained on the Anthropic HH dataset using PPO \footnote{\url{https://huggingface.co/reciprocate/ppo_hh_pythia-6B}} and released by a prominent third party \footnote{\url{https://github.com/CarperAI/trlx/tree/main/examples/hh}}; however, we are unable to identify any prompt or sampling temperature that enables it to outperform the foundational Pythia-2.8B. Drawing on insights from the TL;DR evaluations and recognizing that both Best of 128 and PPO target the same reward objective, we treat Best of 128 as a rough stand-in for PPO-level outcomes. In summary, DPO emerges as the sole approach that is both computationally efficient and capable of exceeding the Anthropic HH dataset’s preferred completions, while offering comparable or superior results to the far more resource-intensive Best of 128. Lastly, Figure~\ref{fig:dialogue-main} demonstrates that DPO attains peak effectiveness in relatively few training steps.

\subsection{Generalization to a new input distribution}

To more rigorously assess how PPO and DPO perform under distribution shifts, we examine the policies derived from our Reddit TL;DR summarization study on a distinct dataset: news articles from the test split of the CNN/DailyMail corpus \citep{nallapati-etal-2016-abstractive}. We employ the optimal sampling temperatures found in the TL;DR experiments (0 and 0.25) for evaluation. The findings, summarized in Table~\ref{tab:ood}, indicate that we utilized the same GPT-4 (C) prompt as in the TL;DR task, with minor wording changes—replacing “forum post” with “news article”—to compute GPT-4 win rates against the datasets' reference summaries. On this alternative data distribution, DPO continues to outperform PPO by a substantial margin. These results offer preliminary support that DPO policies can generalize as effectively as PPO policies, despite DPO lacking access to the supplementary unlabeled Reddit TL;DR prompts that are available to PPO.

\subsection{Validating GPT-4 judgments with human judgments}
\label{sec:human-judgments}
To evaluate the trustworthiness of GPT-4's judgments, we conducted a human annotation study based on the TL;DR summarization experiment outcomes and leveraging two versions of GPT-4 prompts. The \textbf{GPT-4 (S)} (simple) prompt focuses solely on which summary better encapsulates key information from the post. In contrast, the \textbf{GPT-4 (C)} (concise) prompt explicitly asks for the more concise summary, addressing our observation that GPT-4—when using the (S) prompt—tends to prefer summaries that are longer and more repetitive than those favored by human raters. (For full prompt texts, refer to Appendix~\ref{app:prompts}.) Our analysis covers three model comparisons: the top-performing method (DPO, temp. 0.25), the lowest-performing one (PPO, temp. 1.0), and a model with intermediate results (SFT, temp. 0.25), to capture a spectrum of output quality. Each of these is evaluated against PPO using its optimal (greedy) temperature setting. The data reveal that, across both prompt types, GPT-4’s decisions align with those of human annotators at rates comparable to human-human agreement. This finding implies that GPT-4 serves as an acceptable stand-in for human evaluation (noting, however, that we obtained repeated human ratings only for the DPO and PPO-1 pairings, due to limited rater availability). In general, the \textbf{GPT-4 (C)} prompt yields win rates that more closely reflect human preferences; thus, this prompt is adopted in Section~\ref{sec:dpo-real-datasets} for key results. For comprehensive details regarding the annotation process, the web interface for raters, and a roster of participating annotators, please consult Appendix~\ref{app:human-study}.

\section{Discussion}
Preference-based learning stands out as an effective and adaptable strategy for developing language models that are both powerful and aligned with human values. In this work, we presented DPO, a straightforward methodology that enables training language models from preference data without resorting to reinforcement learning approaches. Instead of reformulating preference learning as a conventional reinforcement learning (RL) task to leverage existing RL techniques, DPO establishes a correspondence between the policy of a language model and reward functions. This direct connection allows models to learn human preferences using a simple cross-entropy loss, avoiding both reinforcement learning and loss of expressivity. Remarkably, DPO matches or outperforms established RLHF baselines, such as those utilizing PPO, even in the absence of extensive hyperparameter optimization. As a result, DPO significantly lowers the complexity involved in teaching language models from preference data.

\textbf{Limitations \& Future Work.} Our findings open several avenues for future investigation. One key question concerns the extent to which DPO-trained policies transfer to out-of-distribution examples compared to those trained with explicit reward models. Preliminary experiments indicate that DPO exhibits generalization capabilities comparable to PPO-based methods, yet broader and more systematic evaluations are necessary. Additional areas to explore include the effectiveness of using DPO for self-labeling on unlabeled prompt data. Furthermore, a deeper examination is needed to understand how reward over-optimization might appear under direct preference optimization and whether the modest performance dip observed in Figure~\ref{fig:dialogue-main}-right exemplifies this issue. While our work focuses on models containing up to 6B parameters, a promising next step involves investigating DPO's scalability to models that are considerably larger. In terms of evaluation practices, we observe that GPT-4-derived win rates can be sensitive to prompt design; further research is warranted to optimize automated evaluation protocols. Beyond the training of language models from preferences, DPO has the potential to be applied in a broad array of generative modeling contexts, including domains outside of text.